% TOP_PROBLEM: {"id":30007751,"problem_number":"LOCAL-30007751","title":"Preventing persistent homelessness","category_id":21,"category":{"id":21,"name":"economics","display_name":"Economics","description":"Open economic theory statements, published research agendas, and proposed scoped specifications; question provenance and historical priority are qualified per record.","slug":"economics","order_index":21,"created_at":"2026-10-08T00:00:00Z"},"status":"research_question","external_url":null,"legacy_ids":[],"published":false,"statement":"Choose budget-feasible homelessness services using treatment benefits rather than predicted risk alone, with twenty-four-month housing and health outcomes.","economics":{"original_id":240,"field":"Urban, housing and spatial economics","kind":"design","formalization_basis":"adapted_research_question","source_status":"research_agenda","priority_status":"earliest_found","priority_note":"Earliest explicit statement located in this audit; no exhaustive priority search or first-ever claim. The mathematical specification below is adapted, not quoted from the source.","earliest_verified_reference":{"authors":["William N. Evans","David C. Phillips","Krista Ruffini"],"title":"Reducing and Preventing Homelessness: A Review of the Evidence and Charting a Research Agenda","year":2019,"url":"https://www.povertyactionlab.org/sites/default/files/research-paper/reducing-preventing-homelessness-review-research-agenda.pdf","doi":"10.3386/w26232","locator":"Research agenda, client targeting and matching discussion, printed pp. 45–46","evidence_summary":"The authors call for experiments on targeting and matching services, treatment heterogeneity, and the distinction between predicted risk and treatment benefit."},"current_reference":{"authors":["William N. Evans","David C. Phillips","Krista Ruffini"],"title":"Reducing and Preventing Homelessness: A Review of the Evidence and Charting a Research Agenda","year":2019,"url":"https://www.povertyactionlab.org/sites/default/files/research-paper/reducing-preventing-homelessness-review-research-agenda.pdf","doi":"10.3386/w26232","locator":"Research agenda, client targeting and matching discussion, printed pp. 45–46","evidence_summary":"The authors call for experiments on targeting and matching services, treatment heterogeneity, and the distinction between predicted risk and treatment benefit."},"formalization_note":"The source explicitly identifies the underlying gap or agenda. Population, horizon, interventions, estimands, and auxiliary assumptions are proposed operational choices, not a canonical source theorem. Partial later answers are compatible with this scoped question; the citation alone does not prove absence of every subsequent solution. Mathematical scope was checked independently; added definedness, feasibility, or identification conditions belong to this proposed formalization. Mathematical scope was checked independently; added definedness, feasibility, or identification conditions belong to this proposed formalization."},"statusQualification":"Published question or research agenda verified at the cited locator. The mathematical specification is adapted for this catalogue and includes additional assumptions; source publication does not establish first-ever priority or certify this exact model remains unsolved. The source explicitly identifies the underlying gap or agenda. Population, horizon, interventions, estimands, and auxiliary assumptions are proposed operational choices, not a canonical source theorem. Partial later answers are compatible with this scoped question; the citation alone does not prove absence of every subsequent solution. Mathematical scope was checked independently; added definedness, feasibility, or identification conditions belong to this proposed formalization. Mathematical scope was checked independently; added definedness, feasibility, or identification conditions belong to this proposed formalization.","statusReviewedAt":"2026-10-08"}
% FIRST_POSTED: 2026-10-08
% ENTRY_KIND: statement_only
% !TeX program = lualatex
\documentclass[11pt]{article}
\usepackage{fontspec}
\usepackage[a4paper,margin=25mm]{geometry}
\usepackage{amsmath,amssymb,mathtools}
\usepackage{xurl}
\usepackage[hidelinks]{hyperref}
\newcommand{\sref}[2]{\hyperlink{source:#1:#2}{[S#2]}}
\setlength{\emergencystretch}{3em}
\begin{document}
% problemId:30007751
\section*{Preventing persistent homelessness}\label{30007751}
\subsection{Definitions and mathematical statement}
Consider adults entering a coordinated homelessness-assessment system in one city during a specified year. For person \(i\), let baseline covariates \(X_i\) include housing history, health, income, and family composition. A feasible allocation rule \(d(X_i)\in\{0,1,2\}\) assigns existing services, temporary rental assistance, or permanent supportive housing. Let \(H_i(d)\) be days without stable housing during the next twenty-four months, \(Q_i(d)\) health-related utility, and \(C_i(d)\) real service cost. For a per-eligible-person budget \(B\) and health weight \(\lambda>0\), define
\[d^*\in\arg\max_d E[\lambda Q_i(d(X_i))-H_i(d(X_i))]\quad\text{subject to}\quad E[C_i(d(X_i))]\leq B.\]
Restrict rules to a declared feasible class with finite outcome means. Establish attainment or report the supremum and near-optimal rules. State capacity and equity constraints before optimization. Identify treatment-response heterogeneity through randomized offers or a justified allocation discontinuity, distinguishing predicted homelessness risk from predicted treatment benefit. Follow people across shelters, municipalities, and stable housing exits, with explicit missing-data assumptions. Include crowd-out and housing-market spillovers where program scale makes them relevant. The source explicitly requests targeting and matching experiments and warns that risk and benefit differ. This adapted allocation problem fixes the population, horizon, objectives, and treatments rather than suppose one universal program prevents every recurrence.

\par\textbf{Formulation and scope.} The source explicitly identifies the underlying gap or agenda. Population, horizon, interventions, estimands, and auxiliary assumptions are proposed operational choices, not a canonical source theorem. Partial later answers are compatible with this scoped question; the citation alone does not prove absence of every subsequent solution. Mathematical scope was checked independently; added definedness, feasibility, or identification conditions belong to this proposed formalization. Mathematical scope was checked independently; added definedness, feasibility, or identification conditions belong to this proposed formalization.

\par\textbf{Open-question provenance.} An explicit question or research agenda was verified at the source locator. This does not by itself certify that every later version remains unresolved \sref{30007751}{1}.

\par\textbf{Historical priority.} Earliest explicit statement located in this audit; no exhaustive priority search or first-ever claim. The mathematical specification below is adapted, not quoted from the source.

\subsection{Short English statement}
Choose budget-feasible homelessness services using treatment benefits rather than predicted risk alone, with twenty-four-month housing and health outcomes.

\subsection{Sources}
\begin{itemize}\raggedright
\item[\textnormal{[S1]}]\hypertarget{source:30007751:1}{} William N. Evans, David C. Phillips, Krista Ruffini. 2019. Reducing and Preventing Homelessness: A Review of the Evidence and Charting a Research Agenda. Research agenda, client targeting and matching discussion, printed pp. 45–46.
\url{https://www.povertyactionlab.org/sites/default/files/research-paper/reducing-preventing-homelessness-review-research-agenda.pdf}.
DOI: 10.3386/w26232.
{\small\textit{Earliest verified question/agenda reference; historical priority qualified below. The authors call for experiments on targeting and matching services, treatment heterogeneity, and the distinction between predicted risk and treatment benefit.}}
\end{itemize}
\end{document}
